\documentclass[a4paper,11pt]{article}
\usepackage[T1]{fontenc}
\usepackage{lmodern}
\usepackage[margin=2.5cm]{geometry}
\usepackage[hidelinks]{hyperref}
\input{glyphtounicode}
\pdfgentounicode=1
\pagestyle{empty}
\setlength{\parindent}{0pt}
\setlength{\parskip}{0.9em}
\raggedright

\begin{document}

{\Large\bfseries Abbas Taqvi}\\[4pt]
Milan, Italy \enspace | \href{mailto:abbastaqvi81@gmail.com}{abbastaqvi81@gmail.com}\\
\href{https://github.com/smabbasht}{github.com/smabbasht} \enspace | +39 340 065 5647

\vspace{1.2em}
\textbf{Application: Software Engineering Intern -- Fullstack}

Dear Doctolib Engineering Team,

Publishing Bayaaz, my poetry app, meant building more than a mobile interface. I shipped it on Google Play and the App Store, and I maintain its FastAPI backend, PostgreSQL database and React frontend. It gave me a practical reason to work across the stack: each piece had to come together into something people could use.

At Logiciel, I worked on a different kind of interface: a Python AI chatbot that let business teams query financial data in plain language through a REST API. It reduced analyst turnaround by 41\%. Alongside that, I built backend infrastructure and a Rust framework that cut daily processing from 500 to 50 minutes. These projects taught me to connect technical decisions to a concrete problem for the person using the system.

That is what draws me to Doctolib: building features that make daily work easier for care teams and patients. I am now pursuing an M.S. in Data Science at Milano-Bicocca and would like to bring my experience shipping software to your Milan team. I use AI-assisted coding tools, automated testing and peer review, and would welcome the chance to learn your stack alongside experienced engineers.

Thank you for considering my application.

Best regards,\\
Abbas Taqvi

\end{document}
